% Slide 6 -- Part B4: Comparison of Time-Slicing, NVIDIA MPS, and MIG.
\begin{frame}[t]{So sánh Time-Slicing, NVIDIA MPS và MIG}
\begin{canvas}
% Header row (y = 5.66 .. 6.24)
\fill[BK,rounded corners=2.5pt] (0,5.66) rectangle (2.68,6.24);
\node[anchor=west,text=white,font=\fontsize{9.2}{11}\selectfont\bfseries] at (.16,5.95){Tiêu chí};

\fill[Amber!90,rounded corners=2.5pt] (2.78,5.66) rectangle (6.62,6.24);
\node[anchor=west,text=white,font=\fontsize{9.2}{11}\selectfont\bfseries] at (2.94,5.95){Time-Slicing};

\fill[BK!82,rounded corners=2.5pt] (6.72,5.66) rectangle (10.56,6.24);
\node[anchor=west,text=white,font=\fontsize{9.2}{11}\selectfont\bfseries] at (6.88,5.95){NVIDIA MPS};

\fill[Teal,rounded corners=2.5pt] (10.66,5.66) rectangle (14.55,6.24);
\node[anchor=west,text=white,font=\fontsize{9.2}{11}\selectfont\bfseries] at (10.82,5.95){MIG (Multi-Instance)};

% 5 comparison rows with alternating subtle backgrounds
\foreach \y/\bg/\crit/\cOne/\cTwo/\cThree [count=\r] in {
 4.74/BK!3/Cơ chế chia sẻ/{Luân phiên lát cắt thời gian\\(Context Switching)}/{Chạy song song nhiều tiến\\trình CUDA qua MPS Server}/{Chia phần cứng thành\\tối đa 7 instances độc lập},
 3.82/white/Mức độ cô lập/{Phần mềm; \textbf{không ép VRAM}\\(nguy cơ OOM chéo)}/{Phần mềm (thread \%, VRAM);\\lỗi nghiêm trọng dễ lan}/{Phần cứng hoàn toàn\\(SM, VRAM, băng thông, lỗi)},
 2.90/BK!3/Hiệu năng \& QoS/{Có overhead chuyển ngữ\\cảnh; latency biến động}/{Throughput cao, giảm mạnh\\context switching}/{QoS ổn định, dự đoán được\\không nhiễu giữa các slice},
 1.98/white/GPU \& GKE/{Rộng (Pascal+); tối đa\\48 client/GPU vật lý}/{Volta+ tối ưu ($\le$ 48 client);\\cần \texttt{hostIPC: true}}/{Chỉ Ampere+ (A100, H100);\\tối đa 7 slice/GPU},
 1.06/BK!3/Workload phù hợp/{Dev/test, notebook, tải\\bursty có khoảng nghỉ}/{Batch nhỏ/vừa, tiến trình\\hợp tác (MPI/CUDA)}/{Inference production cần\\SLA và cô lập lỗi chặt}
}{
 \fill[\bg,rounded corners=2pt] (0,\y) rectangle (2.68,\y+.82);
 \draw[BK!18,rounded corners=2pt] (0,\y) rectangle (2.68,\y+.82);
 \node[anchor=west,text=BK,font=\fontsize{8.6}{10.2}\selectfont\bfseries,text width=2.42cm] at (.14,\y+.41){\crit};

 \fill[\bg,rounded corners=2pt] (2.78,\y) rectangle (6.62,\y+.82);
 \draw[BK!18,rounded corners=2pt] (2.78,\y) rectangle (6.62,\y+.82);
 \node[anchor=west,text=Ink,font=\fontsize{8.2}{10}\selectfont,text width=3.58cm] at (2.92,\y+.41){\cOne};

 \fill[\bg,rounded corners=2pt] (6.72,\y) rectangle (10.56,\y+.82);
 \draw[BK!18,rounded corners=2pt] (6.72,\y) rectangle (10.56,\y+.82);
 \node[anchor=west,text=Ink,font=\fontsize{8.2}{10}\selectfont,text width=3.58cm] at (6.86,\y+.41){\cTwo};

 \fill[\bg,rounded corners=2pt] (10.66,\y) rectangle (14.55,\y+.82);
 \draw[BK!18,rounded corners=2pt] (10.66,\y) rectangle (14.55,\y+.82);
 \node[anchor=west,text=Ink,font=\fontsize{8.2}{10}\selectfont,text width=3.62cm] at (10.80,\y+.41){\cThree};
}

\accentbar{Teal}{Kết hợp thực tế}{Trên GKE có thể chạy Time-Slicing hoặc MPS trên từng MIG partition để tăng mật độ container.}
\end{canvas}
\end{frame}
